\section{Conclusions}
\label{sec:conclusions}

The project found the frequent groups of actors among the IMDb Top 1000 movies and compared four methods to find them. At $s = 2$ there are 121 frequent pairs, 25 triples and 4 groups of four, mostly the recurring casts of film series and long-running collaborations; since most actors appear in a single movie, frequent groups are few and small.

All methods return identical results, checked against exhaustive counting in every run and against the known result on the enlarged data. With at most four actors per basket, exhaustive counting is hard to beat on the original data, but its cost grows exponentially with the basket size. Basic A-Priori is dominated by building all $\binom{m}{2}$ pairs of frequent actors as candidates; counting only the pairs that occur, joining within prefix groups and trimming the baskets remove this cost without changing the result, and make optimized A-Priori the fastest method wherever pruning is possible. PCY needs the fewest pair counters in the second pass, but it is consistently slower, its filter loses effectiveness as the average bucket count approaches $s$, and in this implementation it does not reduce total memory.

Under both enlargements, optimized A-Priori and PCY grow approximately linearly with the number of baskets, while basic A-Priori grows roughly quadratically as soon as the actors grow with the data. This confirms that the scalability of A-Priori depends on keeping the candidates and counters proportional to what actually occurs in the data, not only on the number of passes.

\paragraph{Limitations.} The baskets contain at most four actors, which keeps exhaustive counting feasible and the frequent itemsets small. The enlarged datasets are controlled workloads that repeat existing co-occurrences: they show how the costs respond to more baskets and to more distinct actors, not how real data would grow. Everything runs in memory on one machine, so a pass is a loop over a list rather than a read from disk, and memory is described by the number of counters rather than measured in bytes.

\paragraph{Extensions.} Since much of the remaining cost is reading all the baskets once per pass, the limited-pass algorithms of RU \S6.4 are the natural next step. Toivonen's algorithm~\cite{toivonen1996} mines a random sample in memory and verifies the result with one full pass over the data, repeating with a new sample when the verification fails; the SON algorithm~\cite{savasere1995} mines chunks of the data independently and combines them in a second pass, and can be distributed over a cluster with MapReduce or Spark.
